\section{Two boundary experiments at a common throat}
A three-form flux supports a smooth minimum of the area radius of a Euclidean three-sphere \cite{GS}. The geometry remembers the magnitude of that flux. A spinor also sees its sign through the axial component of the spin connection. Does a spinor then distinguish the two branches at the reflection surface? A boundary value at zero eigenvalue needs a domain, a Fock vacuum needs a polarization, and a determinant needs a regulator. We perform two explicit experiments with these choices fixed and compare their dependence on the same three-dimensional data.

We call the portion from the reflection surface to one outer end a \emph{semi-wormhole}, following the geometrical use in \cite{Kaul}. A quotient acts freely on every round section and leaves the radial equation unchanged. The allowed spinor modes can change substantially. Left and right order-three quotients have the same Einstein action and opposite asymmetric spin spectra. The lens space used below gives a different diagnostic: its two lowest contacting shells are absent, but its eta invariant is nonzero. This separates a measurement of occupied states at the throat from a measurement using all eigenvalues.

In the first experiment an outer Euclidean boundary prepares a terminal vacuum bra. A Lorentzian initial surface prepares an instantaneous vacuum ket. Both are transported to the same reflection surface using the chosen coframe. Only a finite-dimensional shell can distinguish the two covariances in the first crossing window. A single Weyl field counts this difference with a sign. In a Dirac field the opposite chirality cancels its vector charge and doubles its axial charge. This distinction survives even when both states have a vanishing vacuum overlap.

In the second experiment both ends carry spectral boundary conditions. The radial Dirac equation separates into pairs of mutually adjoint first-order operators. Their boundary Wronskians yield the same log ratio for every shell of a given sign. Summing that constant requires the signed degeneracies at arbitrarily large angular momentum. We continue the compact four-dimensional zeta difference before differentiating at zero. This identifies the angular factor with the ordinary Dirac eta invariant in the stated orientation. We give an independent shell-first prescription at infinity.

The background and many ingredients of the calculation have precedents. Kaul and Sengupta connected torsional instantons and the Nieh--Yan invariant \cite{Kaul}. Kim studied normalizable fermion zero modes on an Einstein--Kalb--Ramond wormhole \cite{Kim}. Neither its noncompact domain nor the absolute number of such modes specifies the collar experiment. Loges, Shiu and Sudhir found stability in the three-form representation \cite{LSS}. Marolf and Missoni retained the regular symmetric modes when solving the constraints \cite{MM}. We use their bulk scalar quadratic form and trace the endpoint terms through the change of ensemble. Other choices of field variables and boundary conditions have led to different putative negative directions \cite{HTV}. We distinguish a finite outer boundary from the limiting action at infinity.

A separate enhancement mechanism occurs in metric-affine axion gravity \cite{TY}. In their Eqs.~(12)--(14), (23)--(24), the nonminimal functions are functions of the modulus $f$ of a complex scalar. The torsion and nonmetricity coefficients multiply $\partial f$; the phase carries the PQ charge. The equations for the radius and $f$ are coupled, with $f'(0)=0$ and $f(\infty)=f_a$. Freezing $f$ removes that connection source. Here the independent metric-compatible connection couples directly to the derivative of the shift field, and its constant Holst coefficient rescales the charge kinetic term. The square root in Sec.~\ref{sec:background} follows from this one-field normalization. We make no priority claim for an increase in the Euclidean action. Our results concern specified spin lifts, outer domains and preparations.

We use the charge and lift conventions of the companion Paper~08 \cite{P08}. The signs and all angular inputs required here are specified again in Secs.~\ref{sec:background}, \ref{sec:det}. The external massless field, the identification of boundary spinors and the integration contour are physical inputs.
